\section{Composición de movimientos}

El tiro parabólico es la combinación de un MRU en el eje \(x\) y un MRUA en el eje \(y\).

\subsection{Velocidades y posiciones}
\begin{align}
    v_{0x} &= v_0 \cos \theta \\
    v_{0y} &= v_0 \sin \theta \\
    x &= x_0 + v_{0x} \cdot t \\
    y &= y_0 + v_{0y} \cdot t - \frac{1}{2} g t^2
\end{align}

\subsection{Altura máxima}
La altura máxima se alcanza cuando \(v_y = 0\):
\begin{equation}
    y_{\max} = \frac{v_0^2 \sin^2 \theta}{2g}
\end{equation}

\subsection{Alcance horizontal}
El alcance horizontal es la distancia total recorrida en el eje \(x\):
\begin{equation}
    x_{\max} = \frac{v_0^2 \sin(2\theta)}{g}
\end{equation}

\subsection{Tiempo de vuelo}
\begin{equation}
    t_{\text{vuelo}} = \frac{2 v_0 \sin \theta}{g}
\end{equation}

\resumebox{ejercicio}{Fuerza magnética del protón}{\footnotesize
    Una pelota es lanzada con una velocidad de \(20 \, \text{m/s}\) y un ángulo de \(30^\circ\). Calcula:
    \begin{enumerate}
        \item Su altura máxima.
        \item El alcance horizontal.
        \item El tiempo de vuelo.
    \end{enumerate}
}